% !TEX root = ../algebra_02.tex

\section{Свободная группа, аксиоматическое определение}

\begin{Definition}{Свободная группа}{}
	Пусть $F$ — группа, $f_1, \dots, f_n$ — её элементы.
	Группа $F$ называется свободной со свободными образующими $f_1, \dots, f_n$, если выполняется следующее универсальное свойство: для любой группы $G$ и любых элементов $g_1, \dots, g_n \in G$ существует единственный гомоморфизм $\varphi: F \to G$ такой, что $\varphi(f_i) = g_i$ для $i = 1, \dots, n$.
\end{Definition}

\begin{Example}{Свободная группа с одним образующим}{}
	Группа целых чисел $\mathbb{Z}$ является свободной группой со свободным образующим $1$.
	Действительно, для любого $g \in G$ существует гомоморфизм $\varphi: \mathbb{Z} \to G$, заданный правилом $k \mapsto g^k$, при котором $1 \mapsto g$.
\end{Example}

\begin{Proposition}{Единственность свободной группы}{}
	Пусть $F$ и $F'$ — свободные группы со свободными образующими $f_1, \dots, f_n$ и $f_1', \dots, f_n'$ соответственно.
	Тогда существует единственный изоморфизм $\varepsilon: F \to F'$ такой, что $\varepsilon(f_i) = f_i'$ для $i = 1, \dots, n$.
\end{Proposition}

\begin{proof}
	По определению свободной группы существует единственный гомоморфизм $\varepsilon: F \to F'$ с требуемым свойством $\varepsilon(f_i) = f_i'$.
	Аналогично, в обратную сторону существует единственный гомоморфизм $\varepsilon': F' \to F$ такой, что $\varepsilon'(f_i') = f_i$ для $i = 1, \dots, n$.
	Рассмотрим композицию $\varepsilon' \circ \varepsilon: F \to F$. Для образующих имеем:
	\[ (\varepsilon' \circ \varepsilon)(f_i) = f_i, \quad i = 1, \dots, n \]
	Тождественное отображение $id_F: F \to F$ также удовлетворяет условию $id_F(f_i) = f_i$.
	Из единственности гомоморфизма следует, что $\varepsilon' \circ \varepsilon = id_F$.

	Аналогично показывается, что $\varepsilon \circ \varepsilon' = id_{F'}$.
	Поскольку гомоморфизм $\varepsilon$ обратим, он является изоморфизмом.
\end{proof}

\begin{Proposition}{Представление конечно порожденной группы}{}
	Пусть группа $G = \langle g_1, \dots, g_n \rangle$ конечно порождена, а $F$ — свободная группа с $n$ свободными образующими.
	Тогда существует нормальная подгруппа $N \triangleleft F$ такая, что $G \cong F/N$.
\end{Proposition}

\begin{proof}
	Пусть $f_1, \dots, f_n$ — свободные образующие группы $F$. По определению свободной группы существует гомоморфизм $\varphi: F \to G$, переводящий $f_i \mapsto g_i$.
	Так как $g_1, \dots, g_n \in \operatorname{Im} \varphi$, а эти элементы порождают $G$, получаем, что $\langle g_1, \dots, g_n \rangle \subset \operatorname{Im} \varphi$, следовательно, $\operatorname{Im} \varphi = G$.
	По теореме о гомоморфизме:
	\[ \operatorname{Im} \varphi \cong F / \ker \varphi \]
	Поскольку $\operatorname{Im} \varphi = G$, достаточно положить $N = \ker \varphi$.
\end{proof}
